% Part: first-order-logic
% Chapter: proof-systems
% Section: natural-deduction

\documentclass[../../../include/open-logic-section]{subfiles}

\begin{document}

\iftag{FOL}
      {\olfileid{fol}{prf}{ntd}}
      {\olfileid{pl}{prf}{ntd}}

\olsection{नैसर्गिक निगमन}

नैसर्गिक निगमन ही प्रत्यक्ष तर्कविचाराचे अनुकरण करण्यासाठी घडवलेली
!!a{derivation} पद्धत आहे (विशेषतः गणितज्ञ वापरतात त्या नियमबद्ध
तर्कविचाराचे). प्रत्यक्ष तर्कविचार अनेक ``नैसर्गिक'' पद्धतींनी पुढे जातो.
उदाहरणार्थ, प्रकरणांनुसार सिद्धतेमुळे विकल्पयोगी आधारविधानाच्या आधारे
निष्कर्ष प्रस्थापित करता येतो: विकल्पयोगाच्या प्रत्येक घटकापासून
तो निष्कर्ष निष्पन्न होतो, हे स्वतंत्रपणे दाखवले जाते. अप्रत्यक्ष सिद्धतेत
निष्कर्षाचे नकरण व्याघाताकडे नेते हे दाखवून निष्कर्ष प्रस्थापित करता येतो.
सोपाधिक सिद्धतेत पूर्वांगापासून उत्तरांग निष्पन्न होते हे दाखवून
``जर \dots तर \dots'' असा सोपाधिक दावा प्रस्थापित केला जातो. नैसर्गिक
निगमन म्हणजे अशा नैसर्गिक अनुमानांपैकी काहींचे आकारिकीकरण होय. प्रत्येक
तार्किक संयोजक आणि संख्यापक यांना दोन नियम असतात—एक प्रवेशन नियम आणि एक
विलोपन नियम—आणि हे नियम प्रत्येकी अशाच एका नैसर्गिक अनुमान-पद्धतीशी
जुळतात. उदाहरणार्थ, $\Intro{\lif}$ सोपाधिक सिद्धतेशी, तर $\Elim{\lor}$
प्रकरणांनुसार सिद्धतेशी जुळतो. $\Elim{\land}$ हा विशेषतः सोपा नियम आहे;
त्यामुळे $!A \land !B$ पासून~$!A$ (किंवा $!B$) असा निष्कर्ष काढता येतो.

नैसर्गिक निगमनाला इतर !!{derivation} पद्धतींपासून वेगळे करणारे एक
वैशिष्ट्य म्हणजे त्यात गृहीतकांचा होणारा वापर. नैसर्गिक निगमनातील
!!^a{derivation} हा !!{formula}s चा वृक्ष असतो. या !!{formula}s च्या
वृक्षाच्या मुळाशी एकच !!{formula} असते, आणि वृक्षाची ``पाने'' ही अशी
!!{formula}s असतात की ज्यांच्यापासून निष्कर्ष काढला जातो. नैसर्गिक
निगमनात पानांवरील काही !!{formula}s ही !!{derivation} मध्ये भूमिका बजावतात;
पण !!{derivation} निष्कर्षापर्यंत पोहोचेपर्यंत ती ``वापरून संपवलेली'' असतात.
प्रत्यक्ष तर्कविचारात फक्त थोड्या काळापुरती लागू असणारी गृहीतके मांडण्याच्या
प्रथेशी हे जुळते. उदाहरणार्थ, प्रकरणांनुसार सिद्धतेत आपण विकल्पयोगाच्या
प्रत्येक घटकाची सत्यता गृहीत धरतो; सोपाधिक सिद्धतेत पूर्वांगाची सत्यता
गृहीत धरतो; आणि अप्रत्यक्ष सिद्धतेत निष्कर्षाच्या नकरणाची सत्यता गृहीत
धरतो. अशी तात्पुरती गृहीतके मांडणे आणि एखादी मधली पायरी प्रस्थापित करण्यासाठी
ती नंतर बाजूला काढणे, हे नैसर्गिक निगमनाचे खास वैशिष्ट्य आहे. नैसर्गिक
निगमनातील !!{derivation} च्या पानांवरील सूत्रांना गृहीतके म्हणतात,
आणि काही निगमन नियम त्यांना ``!!{discharge}'' करू शकतात. उदाहरणार्थ, काही
गृहीतकांपासून~$!B$ चे असे !!a{derivation} असेल आणि त्या गृहीतकांत~$!A$
असेल, तर $\Intro{\lif}$ नियमामुळे~$!A \lif !B$ असा निष्कर्ष काढता येतो
आणि~$!A$ या रूपातील कोणतेही गृहीतक मुक्त करता येते. (कोणती गृहीतके
कोणत्या अनुमानाच्या वेळी मुक्त केली जातात याची नोंद ठेवण्यासाठी, त्या
अनुमानाला आणि त्याने मुक्त केलेल्या गृहीतकांना
एक क्रमांक देतो.) !!{derivation} च्या शेवटी जी गृहीतके !!{undischarged}
राहतात, ती निष्कर्षाच्या सत्यतेसाठी एकत्रितपणे पुरेशी असतात; म्हणून
!!a{derivation} हे प्रस्थापित करते की त्याची !!{undischarged} गृहीतके
त्याचा निष्कर्ष तार्किकरीत्या निष्पन्न करतात.

नैसर्गिक निगमनावर आधारित $\Gamma \Proves !A$ हा संबंध तेव्हा आणि केवळ
तेव्हाच लागू होतो, जेव्हा असे !!a{derivation} असते की त्यात वृक्षातील
शेवटचे !!{sentence} हे~$!A$ असते आणि !!{undischarged} असलेले प्रत्येक पान
हे~$\Gamma$ मध्ये असते. नैसर्गिक निगमनात $!A$ हे तेव्हा आणि केवळ तेव्हाच
प्रमेय असते, जेव्हा असे !!a{derivation} असते की त्यातील शेवटचे
!!{sentence} हे~$!A$ असते आणि सर्व गृहीतके !!{discharged} झालेली असतात.
उदाहरणार्थ, $\Proves (!A \land !B) \lif !A$ हे दाखवणारे !!a{derivation}
पुढे दिले आहे:
\begin{prooftree}
  \AxiomC{$\Discharge{!A \land !B}{1}$}
  \RightLabel{\Elim{\land}}
  \UnaryInfC{$!A$}
  \DischargeRule{\Intro{\lif}}{1}
  \UnaryInfC{$(!A \land !B) \lif !A$}
\end{prooftree}
$1$ ही खूण दाखवते की $!A \land !B$ हे गृहीतक \Intro{\lif} अनुमानाच्या
वेळी !!{discharged} केले जाते.
  
नैसर्गिक निगमनात संच~$\Gamma$ तेव्हा आणि केवळ तेव्हाच विसंगत असतो,
जेव्हा $\Gamma \Proves \lfalse$. \FalseInt{} नियमामुळे विसंगत संचापासून
कोणतेही !!{sentence} !!{derive}d करता येते.

नैसर्गिक निगमन पद्धती १९३० च्या दशकात गरहार्ड गेंटझेन आणि स्टॅनिस्लॉ
यास्कोव्हस्की (Stanis\l{}aw Ja\'skowski) यांनी विकसित केल्या; पुढे डॅग
प्रावित्झ आणि फ्रेडरिक फिच यांनी त्यांचा आणखी विकास केला. त्यांतील अनुमाने
सिद्धतेच्या नैसर्गिक पद्धतींचे अनुकरण करत असल्यामुळे तत्त्वज्ञ नैसर्गिक
निगमनाला पसंती देतात.
फिच यांनी विकसित केलेल्या आवृत्त्या तर्कशास्त्राच्या प्रास्ताविक
पाठ्यपुस्तकांमध्ये अनेकदा वापरल्या जातात. तर्कशास्त्राच्या तत्त्वज्ञानात
नैसर्गिक निगमनाचे नियम तार्किक कारकांचे अर्थ देतात, असे काही वेळा मानले
गेले आहे (``सिद्धता-उपपत्तीय चिन्हार्थमीमांसा'').

\end{document}
